\section{Project Management and Reflection}\label{sec:reflection}

\subsection{Project structure and milestone management}
\label{sec:reflect-pm}

The project was scoped at the outset as a three-tier validation
pipeline (synthetic, simulated, real-world) extending a single source
paper \cite{morimura2013}. This decomposition gave the work a useful
order: first verify the estimator on the paper's own benchmark, then test
the full detector in labelled simulation, then stress-test the same
pipeline on real data where the labels are only proxies.

Milestone management was driven by a continuously updated
\code{research\_logbook.md}. After each experimental cycle I recorded the
hypothesis, setup, result, interpretation, and next action. The logbook
became especially useful when results changed direction: it made it clear
which claims were still supported by evidence and which were only working
assumptions from an earlier stage.

Two practices worked well:
\begin{itemize}
\item \textbf{Audit-before-extend.} An early retraction of a headline
number made it clear that every extension needed an upstream audit first.
The gradient finite-difference check at $\sim 10^{-9}$ relative error
confirmed the analytic gradient derivation before the feature sweep, and
the contrast-correlation diagnostic showed that AUC alone was not enough
to claim chain recovery.
\item \textbf{Keeping decisions traceable.} Linking each reported number
back to a script, seed, log entry and interpretation note made the late
rewrite possible. It was particularly important for the $f+g$ diagnosis,
the SUMO multistart check and the Labour-Day stress test, all of which
changed how the results should be framed.
\end{itemize}

Two practices that did not work as well are worth recording explicitly:
\begin{itemize}
\item \textbf{Underestimating the compute scaling of the $f+g$ loss.}
The first $f+g$ sweep at SUMO scale was projected at $\sim 1$ hour and
killed at $9.5$ hours. The retry at smaller $|X_o|$ and
$\code{maxiter}$ was projected at $\sim 30$ minutes and killed at
$3$~hours. Two iterations were burned before the structural
diagnosis (the hitting-rate gradient is $O(|X_o| \cdot n^3)$ and the
$84\%$-floor loss landscape is degenerate at this scale) was reached.
In future I would profile a single iteration before committing to a
sweep at unfamiliar configuration.
\item \textbf{Late identification of the OD-mix confound.} The Tier-3
rush-vs-off-peak proxy was scoped before the methodological cost of
the OD-mix confound was clearly understood. The OD-rebalancing
extension, the methodologically correct fix, was added late in the
project, and the matched-pool result it produced
(\cref{sec:phase5-interpretation}) materially changed the Tier-3 interpretation.
Earlier scoping would have caught this; the lesson is that
``proxy weakness'' should be a structured pre-experiment audit item,
not an ad-hoc discussion item.
\end{itemize}


Looking back, the strongest parts of the project are the staged
validation structure and the willingness to keep negative results in the
main argument. The work did not stop at reproducing the source paper: it
tested the estimator inside a detector, separated labelled simulation from
proxy-labelled real data, and used diagnostics to avoid over-reading a
single AUC. The main weaknesses are also clear. Tier~3 uses one real
city/month pool rather than multiple independent real-world replications,
and the feature result that first looked like a detection improvement
became seed-dependent once the full SUMO sweep was run.

\subsection{Personal skill development}
\label{sec:reflect-skills}

The technical skills I leaned on heavily and improved in: deriving and
finite-difference-verifying analytic gradients (Eqs.~12, 15, 17--18 of
the source paper); writing software with a clean separation between
the inverter's mathematical core (in \code{morimura.py}, kept faithful
to the paper) and the application-specific loaders that surround it
(\code{sumo\_to\_phase3.py}, \code{real\_data/load\_xuancheng.py});
and writing a logbook that records both the successful and the
unsuccessful experimental cycles in enough detail that I could
critique my own past decisions a month later.

The non-technical skill I would name as the one that most shaped the
project's outcome is \emph{disciplined treatment of negative results}. The
project did not produce the headline ``real-world detector with a
high AUC and clean chain recovery'' that the SUMO Tier-2 result would
have predicted. Instead it produced: a controlled-validation result
($\sim 0.71$ AUC, $69\%$ gap closure on labelled SUMO), a model-class
diagnostic (the empirical $2$-step detector improves the $1$-step
reference, strengthening the case for a parametric $2$-step
extension), and a real-world result ($\sim 0.53$ AUC on a
behavioural proxy without ground-truth labels) that quantifies the gap
between simulated and real-world detectability. Each of these is a
methodologically useful result, in part because they were framed
accurately. The temptation late in the project to treat a stronger
holiday proxy as a headline rescue had to be handled carefully rather
than either embraced uncritically or dismissed. Running it with the same
OD controls, bootstrap intervals, per-row diagnostics and multistart
audit showed the useful version of the negative result: Labour Day has a
real full-observation within-OD signal, but the fitted detector still
falls to chance under validation-selected multistart. Accepting that
result, rather than chasing a higher headline number, is the skill I most
want to carry forward.
